\chapter{Conclusion}

In this thesis, we have investigated several fundamental problems in quantum information theory, primarily focusing on the interplay between semidefinite programming and dimensionality reduction. 
By exploiting the inherent structure of the problem, we have developed algorithms and theoretical frameworks that address the computational bottlenecks associated with the high-dimensional spaces.

The first part of this work addressed the separability problem. 
By introducing a hybrid quantum-classical heuristic, we demonstrated that the effective dimension of the optimization problem can be significantly reduced while maintaining feasibility guarantees. 
Our results for the one-dimensional Ising Hamiltonian on up to 28 qubits illustrate that this approach can effectively approximate the separable ground energy, providing a scalable tool for characterizing entanglement in many-body systems.

In the second part, we extended the concept of dimension reduction to quantum state discrimination. 
We showed that the search for optimal measurements can be reduced from a problem scaling with the dimension $d$ of the space to one scaling with the number of candidate states $N$. 
This $N L$ reduction framework is remarkably general, accommodating various figures of merit including minimum-error, unambiguous discrimination, and the state exclusion problem. 
We first characterize the optimal identifications for quantum changepoint problems under several reward structures, including multiple-changepoint settings that were previously computationally inaccessible. 
Second, we consider a quantum error classification problem and show how our reduction makes it tractable for systems of hundreds
of qubits.

Finally, we explored structural measurement strategies for the state exclusion problem. 
We introduced the pretty bad measurement (PBM) as a natural analogue to the pretty good measurement, proving its efficacy relative to blind guessing. 
Furthermore, our characterization of Boolean property learning from quantum sequences established the global optimality of memoryless ``greedy'' strategies for affine functions. 
This provides a clear boundary for when local measurements suffice and when joint measurements are fundamentally required to extract coarse-grained information from quantum data.

In summary, the results presented here suggest that many ``hard'' problems in quantum information possess a hidden structure that can be exploited for efficient computation. 
As quantum hardware continues to evolve, the hybrid workflows and dimension-reduction techniques developed in this thesis will likely play a crucial role in bridging the gap between theoretical protocols and practical implementation on near-term devices.